
% pnu_sampleKor.tex v1.0
% 본 양식은 서울대학교 전기컴퓨터공학부 학위논문 양식을 활용하여, 부산대학교의 양식에 맞춘 것입니다.
% https://ee.snu.ac.kr/community/notice/academic?bm=v&bbsidx=48811

% 논문작성을 위해서 한글패키지가 설치되어 있어야 합니다.
% MikTeX 을 사용할 것을 권장하며, MikTeX의 설치 및 한글 패키지와 관련된 사항은
% KTUG 한글 TeX 사용자 그룹 http://www.ktug.or.kr을 참조하시길 바랍니다.

% pnuthesis.cls file을 사용합니다.
% 박사의 경우 \documentclass[doctor]{pnuthesis} (default)
% 석사의 경우 \documentclass[master]{pnuthesis} 로 변경하면 됩니다.
% pnuthesis.cls file명이 바뀔경우에는 {pnuthesis} 대신에 {바뀐 file명}을 넣으면 됩니다.
% 주어진 .cls file의 내용을 변경하는 것은 규격에 어긋난 결과를 낼 수 있습니다.
%
% 주어진 .cls file은 부산대학교 논문 규격에 맞춰서 작성되어 있습니다.
% 변경 사항이 있을 경우, 적용하시길 바랍니다.

% 본 양식은 부산대학교 대학원에서 공식적으로 갱신하는 양식이 아닙니다.
% 활용에 따른 책임은 사용자에 있습니다.
% 조판 후 작성 된 것이 *규정에 맞는지* 꼭 확인하세요.

% 석사는 master, 박사는 doctor
% 영문은 english, 국문은 korean
\documentclass[master, korean]{pnuthesis}

% UsePackages
\usepackage{lineno}
\usepackage{indentfirst}  % 섹션/챕터 다음 첫 문단 들여쓰기
% 알고리즘 패키지
\usepackage{algorithm}
\usepackage{algpseudocode}
\usepackage{amsmath} % \operatorname* 사용을 위해 필요
% 그래프 삽입 패키지
\usepackage{graphicx}
\usepackage{subcaption}
\usepackage{mwe}
\usepackage{svg}
% 테이블 패키지
\usepackage[utf8]{inputenc}
\usepackage{kotex}      % 한글 사용을 위한 필수 패키지
\usepackage{booktabs}   % \toprule, \midrule, \bottomrule 사용을 위한 패키지
\usepackage{caption}    % 캡션 설정을 위한 패키지 (선택 사항)
\usepackage{booktabs}
\usepackage{multirow}
\usepackage{threeparttable}  % tablenotes 환경을 위한 패키지
\usepackage{ragged2e}
% 논문 작성을 위한 사전 준비과정
% 논문제목 (국문, 영문), 저자, 제출일, 심사일, 졸업일의 정보를 넣습니다.

% 논문제목을 넣습니다.
% 필히 한글제목과 영문제목 모두 넣어야 합니다.
\title[korean]{국문 논문 제목 (예시)}



\title[english]{English Thesis Title (Example)}

% 저자 정보를 넣습니다.
% 국문성명, 영문성명 모두 넣어야 하며, 특히 국문성명의 경우는 글자사이에 space가 있는 것과 없는 것
% 두 가지 모두를 집어넣어줘야 합니다.
\author[korean]{박 동 진}
\author[english]{Dongjin Park}
\author[nospace]{박동진}

% 지도교수님의 성함을 국문 (영문 논문일 경우 영어로) 넣습니다.
\adviser[korean]{김 태 운}
\adviser[english]{Taewoon Kim}
\adviser[nospace]{김태운}

% 논문 제출일, 논문 심사일을 한글(영문 논문일 경우 영어)로 넣습니다.
\examinationdate[korean]{2025년 12월 18일}
\examinationdate[english]{17, November, 2025}

% 졸업일을 영문식, 한글식 두 가지 방법 모두 넣습니다.
\gradyear[korean]{2026년 2월}
\gradyear[english]{FEBRUARY 2026}


\abstracts[korean]{
    	이곳에 국문 초록을 작성합니다. 연구의 목적, 연구 방법, 주요 결과, 결론 및 기대 효과 등을 간략하게 요약하여 기술하십시오.
	본 템플릿은 부산대학교 대학원 학위논문 작성을 위한 예시 양식입니다. \LaTeX을 처음 사용하는 분들을 위해 기본적인 문서 구조와 그림, 표, 수식, 참고문헌 등을 삽입하는 방법을 포함하고 있습니다.
}

\abstracts[english]{
	Write your English abstract here. Summarize the purpose, methods, main results, and conclusions of your research.
	This template is provided for writing a master's or doctoral thesis at Pusan National University. It demonstrates how to use basic elements such as figures, tables, equations, and references using \LaTeX.
}
%

% 문서의 시작
%
% 위의 정보들을 빠짐없이 채워넣고 document를 시작하면
% 외표지, 내표지(외표지와 동일), 인준지가 자동으로 생성됩니다.

% \linenumbers

\begin{document}
\renewcommand{\baselinestretch}{1.5}    % 본문의 줄간격 조정, 고치거나 삭제하지 마십시오.
\selectfont                             %

\changepage{5mm}{}{}{}{}{}{}{}{-5mm}    %%페이지 여백 재설정. 절대 고치거나 삭제하지 마십시오.

\makelists % 목차생성

% 본문의 시작
% chapter, section의 추가,변경등 모두를 자유롭게 할 수 있습니다.

% 용어 매크로

\newcommand{\IntegratedView}{통합 뷰}
\newcommand{\Registration}{정합}
\newcommand{\ProcessingNode}{처리 노드}
\newcommand{\UAVfull}{무인항공기 (Unmanned Aerial Vehicle, UAV; 이하 ‘드론’)}
\newcommand{\UAV}{드론}

\chapter{서론}
\section{연구 배경}
이곳에 연구 배경을 기술합니다. 본문 내용은 직접 작성하거나, 필요한 경우 외부 파일을 \texttt{\textbackslash input} 하여 사용할 수 있습니다.

참고문헌 인용 예시입니다 \cite{motlagh2023unmanned}. 두 개 이상의 문헌을 인용할 경우 \cite{motlagh2023unmanned, chen2022wildland}와 같이 작성합니다.

\section{연구 목적}
본 연구의 목적을 기술합니다. 해결하고자 하는 문제점과 연구의 방향성을 명확히 제시합니다.

\chapter{본론 작성을 위한 예시}
이 장에서는 \LaTeX 을 사용하여 논문을 작성할 때 자주 사용되는 요소들의 예시를 보여줍니다.

\section{수식 (Equation)}
수식은 \texttt{equation} 환경을 사용하여 작성합니다.
\begin{equation}
    E = mc^2 \label{eq:einstein}
\end{equation}
식 \ref{eq:einstein}은 아인슈타인의 질량-에너지 등가 원리를 나타냅니다.

복잡한 적분 수식의 예시입니다.
\begin{equation}
    f(x) = \int_{-\infty}^{\infty} \hat{f}(\xi)\,e^{2\pi i \xi x} \,d\xi
\end{equation}

\section{그림 (Figure)}
그림은 \texttt{figure} 환경을 사용합니다. \texttt{graphicx} 패키지가 필요하며, 본 템플릿에서는 예시를 위해 \texttt{mwe} 패키지의 이미지를 사용했습니다.

\begin{figure}[htbp]
    \centering
    \includegraphics[width=0.5\textwidth]{example-image-a}
    \caption{기본 그림 예시}
    \label{fig:sample_a}
\end{figure}

그림 \ref{fig:sample_a}은 기본적인 그림 삽입 예시입니다.

\subsection{다중 그림 (Subfigure)}
여러 그림을 나란히 배치할 때는 \texttt{subcaption} 패키지를 활용합니다.

\begin{figure}[htbp]
    \centering
    \begin{subfigure}[b]{0.45\textwidth}
        \centering
        \includegraphics[width=\textwidth]{example-image-b}
        \caption{왼쪽 그림 예시}
        \label{fig:sample_b}
    \end{subfigure}
    \hfill
    \begin{subfigure}[b]{0.45\textwidth}
        \centering
        \includegraphics[width=\textwidth]{example-image-c}
        \caption{오른쪽 그림 예시}
        \label{fig:sample_c}
    \end{subfigure}
    \caption{서브 피규어 예시}
    \label{fig:subfigures}
\end{figure}

그림 \ref{fig:subfigures}는 두 개의 그림을 가로로 배치한 예시입니다.

\section{표 (Table)}
표는 \texttt{table} 환경 내에 \texttt{tabular}를 사용하여 작성합니다. \texttt{booktabs} 패키지를 사용하면 전문적인 품질의 표를 만들 수 있습니다.

\begin{table}[htbp]
    \centering
    \caption{기본 표 예시}
    \label{tab:sample_table}
    \begin{tabular}{lcc}
        \toprule
        \textbf{항목} & \textbf{변수 A} & \textbf{변수 B} \\
        \midrule
        데이터 1 & 10.5 & 23.4 \\
        데이터 2 & 15.2 & 18.7 \\
        데이터 3 & 12.0 & 20.1 \\
        \bottomrule
    \end{tabular}
\end{table}

표 \ref{tab:sample_table}은 기본적인 표의 형태를 보여줍니다.

\section{알고리즘 (Algorithm)}
알고리즘은 \texttt{algorithm} 및 \texttt{algpseudocode} 패키지를 사용하여 작성합니다.

\begin{algorithm}
    \caption{유클리드 호제법 (Euclidean Algorithm)}
    \label{alg:euclid}
    \begin{algorithmic}[1]
        \Require $a, b$: non-negative integers
        \Ensure $\gcd(a, b)$: greatest common divisor of $a$ and $b$
        \State $r \gets a \bmod b$
        \While{$r \neq 0$}
            \State $a \gets b$
            \State $b \gets r$
            \State $r \gets a \bmod b$
        \EndWhile
        \State \Return $b$
    \end{algorithmic}
\end{algorithm}

알고리즘 \ref{alg:euclid}는 최대공약수를 구하는 과정을 나타냅니다.

\section{목록 (Lists)}
항목을 나열할 때는 \texttt{itemize}와 \texttt{enumerate} 환경을 사용합니다.

순서가 없는 목록:
\begin{itemize}
    \item 첫 번째 항목
    \item 두 번째 항목
    \item 세 번째 항목
\end{itemize}

순서가 있는 목록:
\begin{enumerate}
    \item 첫 번째 단계
    \item 두 번째 단계
    \item 세 번째 단계
\end{enumerate}

\section{텍스트 서식 (Text Formatting)}
다양한 텍스트 강조 및 서식을 사용할 수 있습니다.
\begin{itemize}
    \item \textbf{굵은 글씨 (Bold)}: \texttt{\textbackslash textbf\{...\}}
    \item \textit{이탤릭체 (Italic)}: \texttt{\textbackslash textit\{...\}}
    \item \underline{밑줄 (Underline)}: \texttt{\textbackslash underline\{...\}}
    \item \texttt{타자기 글꼴 (Typewriter)}: \texttt{\textbackslash texttt\{...\}}
\end{itemize}

\section{주석 및 인용 (Footnotes and Quotes)}
각주는 본문 내용에 대한 부연 설명이 필요할 때 사용합니다.\footnote{이것은 각주 예시입니다. 페이지 하단에 표시됩니다.}

인용문은 \texttt{quote} 환경을 사용합니다.
\begin{quote}
    ``지식보다 중요한 것은 상상력이다. 지식은 한계가 있지만, 상상력은 세상의 모든 것을 끌어안기 때문이다.'' \\
    - 알베르트 아인슈타인
\end{quote}

\section{코드 삽입 (Code Listing)}
간단한 소스 코드는 \texttt{verbatim} 환경을 사용하여 그대로 출력할 수 있습니다.

\begin{verbatim}
# Python 코드 예시
def hello_world():
    print("Hello, World!")
\end{verbatim}

\section{한국어 더미 텍스트 (Korean Lorem Ipsum)}
긴 문단을 작성할 때의 줄바꿈과 들여쓰기를 확인하기 위한 한국어 더미 텍스트 예시입니다.

모든 국민은 언론·출판의 자유와 집회·결사의 자유를 가진다. 언론·출판에 대한 허가나 검열과 집회·결사에 대한 허가는 인정되지 아니한다. 통신·방송의 시설기준과 신문의 기능을 보장하기 위하여 필요한 사항은 법률로 정한다. 언론·출판은 타인의 명예나 권리 또는 공중도덕이나 사회윤리를 침해하여서는 아니된다. 언론·출판이 타인의 명예나 권리를 침해한 때에는 피해자는 이에 대한 피해의 배상을 청구할 수 있다.

법관은 헌법과 법률에 의하여 그 양심에 따라 독립하여 심판한다. 대법원에 대법관을 둔다. 헌법재판소 재판관의 임기는 6년으로 하며, 법률이 정하는 바에 의하여 연임할 수 있다.

대통령은 제1항과 제2항의 처분 또는 명령을 한 때에는 지체없이 국회에 보고하여 그 승인을 얻어야 한다. 승인을 얻지 못한 때에는 그 처분 또는 명령은 그때부터 효력을 상실한다. 이 경우 그 명령에 의하여 개정 또는 폐지되었던 법률은 그 명령이 승인을 얻지 못한 때부터 당연히 효력을 회복한다. 대통령은 이를 지체없이 공포하여야 한다.

국가는 균형있는 국민경제의 성장 및 안정과 적정한 소득의 분배를 유지하고, 시장의 지배와 경제력의 남용을 방지하며, 경제주체간의 조화를 통한 경제의 민주화를 위하여 경제에 관한 규제와 조정을 할 수 있다. 광물 기타 중요한 지하자원·수산자원·수력과 경제상 이용할 수 있는 자연력은 법률이 정하는 바에 의하여 일정한 기간 그 채취·개발 또는 이용을 특허할 수 있다.

이 헌법중 공무원의 임기 또는 중임제한에 관한 규정은 이 헌법에 의하여 그 공무원이 최초로 선출 또는 임명된 때로부터 적용한다. 제1항의 규정은 이 헌법 시행 당시의 공무원에게는 적용하지 아니한다. 법령의 제정·개정 또는 폐지는 이 헌법에 위배되지 아니하는 범위안에서 효력을 가진다.

모든 국민은 인간으로서의 존엄과 가치를 가지며, 행복을 추구할 권리를 가진다. 국가는 개인이 가지는 불가침의 기본적 인권을 확인하고 이를 보장할 의무를 진다. 모든 국민은 법 앞에 평등하다. 누구든지 성별·종교 또는 사회적 신분에 의하여 정치적·경제적·사회적·문화적 생활의 모든 영역에 있어서 차별을 받지 아니한다.

대한민국은 국제평화의 유지에 노력하고 침략적 전쟁을 부인한다. 국군은 국가의 안전보장과 국토방위의 신성한 의무를 수행함을 사명으로 하며, 그 정치적 중립성은 준수된다. 헌법개정안은 국회가 의결한 후 30일 이내에 국민투표에 부쳐 국회의원선거권자 과반수의 투표와 투표자 과반수의 찬성을 얻어야 한다. 헌법개정안이 제2항의 찬성을 얻은 때에는 헌법개정은 확정되며, 대통령은 즉시 이를 공포하여야 한다.

선거에 있어서 최고득표자가 2인 이상인 때에는 국회의 재적의원 과반수가 출석한 공개회의에서 다수표를 얻은 자를 당선자로 한다. 대통령후보자가 1인일 때에는 그 득표수가 선거권자 총수의 3분의 1 이상이 아니면 대통령으로 당선될 수 없다. 대통령이 궐위된 때 또는 대통령 당선자가 사망하거나 판결 기타의 사유로 그 자격을 상실한 때에는 60일 이내에 후임자를 선거한다.

\chapter{결론}
본 연구의 결론을 작성합니다. 연구 결과의 요약 및 의의, 향후 연구 과제 등을 기술합니다.

\bibliographystyle{ieeetr}
\bibliography{references}

\pagebreak
\acknowledgement
이곳에 감사의 글을 작성합니다. (선택사항)

교수님과 연구실 동료들에게 감사를 표합니다.


\end{document}